\documentclass[a4paper,12pt]{article}

\usepackage[utf8]{inputenc}
\usepackage[T1]{fontenc}
\usepackage[table]{xcolor}
\usepackage{geometry}
\usepackage{longtable}
\usepackage{enumitem}
\usepackage{hyperref}
\usepackage{listings}
\usepackage{titlesec}
\usepackage{booktabs}
\usepackage{array}

\geometry{margin=2.5cm}

\definecolor{criticalcolor}{HTML}{DC3545}
\definecolor{highcolor}{HTML}{FD7E14}
\definecolor{mediumcolor}{HTML}{FFC107}
\definecolor{lowcolor}{HTML}{28A745}

\newcommand{\severitybox}[2]{%
  \colorbox{#1}{\makebox[2.5cm][c]{\textcolor{white}{\textbf{#2}}}}%
}

\lstset{
  basicstyle=\small\ttfamily,
  breaklines=true,
  frame=single,
  numbers=none,
  tabsize=2,
  backgroundcolor=\color{gray!10},
  extendedchars=true,
  inputencoding=utf8,
  literate={—}{---}1
}

\title{\textbf{Performance Report \#06}\\[4pt]
\large Vulkan Engine — Bug Analysis, Performance Problems \& Solutions}
\author{Henrique Rocha}
\date{\today}

\begin{document}
\maketitle

\begin{abstract}
This report presents a comprehensive analysis of the Vulkan engine codebase, covering 26 identified
issues spanning critical bugs, performance anti-patterns, modern Vulkan compliance gaps,
and synchronization hazards. Each issue includes a severity rating, root cause analysis, and
a concrete remediation strategy. An overall classification and remediation priority matrix
are provided at the end.
\end{abstract}

\tableofcontents
\newpage

\section{Scope \& Methodology}

The analysis covered the full codebase under \texttt{vulkan/}, \texttt{main.cpp},
\texttt{StagingRingBuffer.cpp}, and \texttt{IndirectRenderer.cpp}. Each issue was verified
by reading source code and cross-referencing with the Vulkan specification (1.3 \& 1.4).
Findings are categorized as:

\begin{itemize}[noitemsep]
  \item[\textbf{B}] \textbf{Bugs} — functional defects, undefined behavior, or validation errors.
  \item[\textbf{P}] \textbf{Performance} — suboptimal patterns causing CPU/GPU overhead.
  \item[\textbf{C}] \textbf{Compliance} — deviations from modern Vulkan best practices.
  \item[\textbf{S}] \textbf{Synchronization} — incorrect or missing barriers, races, stalls.
\end{itemize}

\section{Issue Register}

\subsection{Critical Severity}

\begin{enumerate}[leftmargin=*]

\item[\textbf{[B1]}] \textbf{Async Command Buffer Ring Unsafe Wrap}
  \hfill \severitybox{criticalcolor}{CRITICAL}

  \begin{tabular}{ll}
    \textbf{File:} & \texttt{vulkan/VulkanApp.cpp} \\
    \textbf{Line:} & 2616--2651 \\
    \textbf{Category:} & Bug / Memory Safety \\
  \end{tabular}

  \textbf{Description:}
  The async command buffer ring (3 graphics + 3 async queues) maintains a fixed-size pool of
  pre-allocated command buffers. When \texttt{asyncSubmitIndex} wraps around after 64
  submissions, it may reuse a command buffer that the GPU is still consuming. There is no
  fence or semaphore wait before recycling.

  \textbf{Root Cause:}
  The ring index is incremented unconditionally without checking whether the previous
  submission using that slot has completed:

\begin{lstlisting}
// Line ~2630 -- simplified
uint32_t slot = asyncSubmitIndex % ASYNC_RING_SIZE;
// No wait for slot availability
vkResetCommandBuffer(cmdBuffers[slot], 0);
// Record and submit
asyncSubmitIndex++;
\end{lstlisting}

  \textbf{Fix:}
  Associate a \texttt{VkFence} (or timeline semaphore value) with each ring slot.
  Before resetting a slot, wait on its fence. This guarantees the GPU has finished
  consuming the previous command buffer.

  \textbf{Severity Justification:} Memory corruption and GPU hangs are guaranteed
  under sustained async workload. Exploitable at runtime.

\item[\textbf{[B2]}] \textbf{Device-Lost Path Skips Cleanup}
  \hfill \severitybox{criticalcolor}{CRITICAL}

  \begin{tabular}{ll}
    \textbf{File:} & \texttt{vulkan/VulkanApp.cpp} \\
    \textbf{Line:} & 931--938 \\
    \textbf{Category:} & Bug / Resource Leak \\
  \end{tabular}

  \textbf{Description:}
  When a \texttt{VK\_ERROR\_DEVICE\_LOST} is caught, the exception handler calls
  \texttt{throw} without invoking \texttt{cleanup()}. This leaks the GLFW window,
  Vulkan instance, device, all allocations, and leaves the GPU in an undefined state.

  \textbf{Root Cause:}
  The exception is caught but not properly handled:

\begin{lstlisting}
// Line ~935
catch (const std::runtime_error& e) {
    std::cerr << "Fatal: " << e.what() << std::endl;
    throw;  // Skips cleanup()
}
\end{lstlisting}

  \textbf{Fix:}
  Replace the bare \texttt{throw} with a call to \texttt{cleanup()} followed by
  \texttt{std::exit(EXIT\_FAILURE)} (or re-throw after cleanup). In the RAII
  pattern, wrap device/instance in unique\_ptr with custom deleters.

  \textbf{Severity Justification:} Resource leak on a fatal path. Prevents
  driver recovery and complicates debugging.

\end{enumerate}

\subsection{High Severity}

\begin{enumerate}[leftmargin=*]

\item[\textbf{[B3]}] \textbf{GPU Buffer Leak on Exception in Async Upload}
  \hfill \severitybox{highcolor}{HIGH}

  \begin{tabular}{ll}
    \textbf{File:} & \texttt{vulkan/VulkanApp.cpp} \\
    \textbf{Line:} & 81--101 \\
    \textbf{Category:} & Bug / Memory Leak \\
  \end{tabular}

  \textbf{Description:}
  The async upload path allocates a staging and a GPU buffer before queue submission.
  If the queue submission or fence creation fails, the allocated buffers are leaked
  because no cleanup path exists between allocation and the exception.

  \textbf{Fix:}
  Wrap buffer allocations in RAII guards (e.g., \texttt{std::unique\_ptr} with a
  custom deleter that calls \texttt{vkDestroyBuffer} and \texttt{vkFreeMemory}).
  Alternatively, allocate via VMA which handles rollback.

\item[\textbf{[S1]}] \textbf{Shadow Cascade Depth Barrier — Wrong Source Stage}
  \hfill \severitybox{highcolor}{HIGH}

  \begin{tabular}{ll}
    \textbf{File:} & \texttt{vulkan/renderer/ShadowRenderer.cpp} \\
    \textbf{Line:} & 354 \\
    \textbf{Category:} & Synchronization / Race Hazard \\
  \end{tabular}

  \textbf{Description:}
  The depth attachment barrier for shadow cascades transitions the depth image from
  \texttt{UNDEFINED} to \texttt{DEPTH\_ATTACHMENT\_OPTIMAL} using
  \texttt{VK\_PIPELINE\_STAGE\_2\_EARLY\_FRAGMENT\_TESTS\_BIT} as the destination
  stage. However, the preceding pass (writing shadow map from a different cascade)
  may still be using \texttt{LATE\_FRAGMENT\_TESTS\_BIT}. The barrier does not wait
  on the late tests stage, creating a hazard.

  \textbf{Fix:}
  Use \texttt{VK\_PIPELINE\_STAGE\_2\_LATE\_FRAGMENT\_TESTS\_BIT} or the union
  \texttt{EARLY\_TESTS\_BIT | LATE\_TESTS\_BIT} as the source stage to ensure all
  previous depth writes complete before the image is reused.

\item[\textbf{[S2]}] \textbf{ALL\_COMMANDS\_BIT Causes Full Pipeline Flush Every Frame}
  \hfill \severitybox{highcolor}{HIGH}

  \begin{tabular}{ll}
    \textbf{File:} & \texttt{main.cpp} \\
    \textbf{Line:} & 872, 895 \\
    \textbf{Category:} & Synchronization / Performance \\
  \end{tabular}

  \textbf{Description:}
  The cubemap rendering and depth prepass barriers both use
  \texttt{VK\_PIPELINE\_STAGE\_2\_ALL\_COMMANDS\_BIT}. This is a
  \textit{full pipeline flush} — it injects a global synchronization point
  that prevents all overlapping execution on the GPU, destroying any
  concurrency between graphics and compute work.

  \textbf{Fix:}
  Replace \texttt{ALL\_COMMANDS\_BIT} with the narrowest stage mask required:
  \texttt{ALL\_GRAPHICS\_BIT}, or better yet the specific stage (e.g.,
  \texttt{COLOR\_ATTACHMENT\_OUTPUT\_BIT} for output reads,
  \texttt{VERTEX\_SHADER\_BIT} for input reads). The principle: synchronize
  only what must be visible, not everything.

\item[\textbf{[P1]}] \textbf{Staging Ring Buffer Lacks Double-Free Protection}
  \hfill \severitybox{highcolor}{HIGH}

  \begin{tabular}{ll}
    \textbf{File:} & \texttt{StagingRingBuffer.cpp} \\
    \textbf{Line:} & 81--131 \\
    \textbf{Category:} & Performance / Correctness \\
  \end{tabular}

  \textbf{Description:}
  The \texttt{release()} method on staging allocations decrements an internal
  reference counter and frees memory when it reaches zero. However, there is
  no guard against double-release of the same allocation, which would cause
  use-after-free of the staging memory and corrupt the free list.

  \textbf{Fix:}
  Set a \texttt{released} flag on the allocation metadata when freed.
  Assert or silently ignore on double-release. Additionally, consider
  using \texttt{std::shared\_ptr} with a custom deleter for reference
  counting instead of a manual counter.

\end{enumerate}

\subsection{Medium Severity}

\begin{enumerate}[leftmargin=*]

\item[\textbf{[B4]}] \textbf{Long Wait on AcquireNextImage — No Timeout Handling}
  \hfill \severitybox{mediumcolor}{MEDIUM}

  \begin{tabular}{ll}
    \textbf{File:} & \texttt{vulkan/VulkanApp.cpp} \\
    \textbf{Line:} & 4800--4810 \\
    \textbf{Category:} & Bug / Robustness \\
  \end{tabular}

  \textbf{Description:}
  \texttt{vkAcquireNextImageKHR} is called with \texttt{UINT64\_MAX} timeout.
  If the window is minimized or the driver hangs, this blocks the main thread
  indefinitely. No fallback or timeout handling is present.

  \textbf{Fix:}
  Reduce the timeout (e.g., 1 second) and handle
  \texttt{VK\_TIMEOUT} by yielding and retrying. For minimized windows,
  throttle to 1~FPS or suspend rendering entirely.

\item[\textbf{[B5]}] \textbf{Heap-Allocated Descriptor Structs Not Exception-Safe}
  \hfill \severitybox{mediumcolor}{MEDIUM}

  \begin{tabular}{ll}
    \textbf{File:} & \texttt{vulkan/renderer/SceneRenderer.cpp} \\
    \textbf{Line:} & 619--796 \\
    \textbf{Category:} & Bug / Exception Safety \\
  \end{tabular}

  \textbf{Description:}
  Eight \texttt{VkDescriptorImageInfo} and \texttt{VkDescriptorBufferInfo}
  structs are allocated with \texttt{new} during initialization then
  \texttt{delete}d later (lines 764--773). If an exception is thrown between
  allocation and deletion, the memory is leaked. Additionally, these are
  small structs that should be stack-allocated.

  \textbf{Fix:}
  Use stack allocation (\texttt{VkDescriptorImageInfo info\{\};}) for all
  descriptor info structs. They are trivially copyable and only need to
  live until \texttt{vkUpdateDescriptorSets} returns.

\item[\textbf{[P2]}] \textbf{Direct vkAllocateMemory Bypasses VMA}
  \hfill \severitybox{mediumcolor}{MEDIUM}

  \begin{tabular}{ll}
    \textbf{File:} & \texttt{vulkan/VulkanApp.cpp} \\
    \textbf{Line:} & 4604 \\
    \textbf{Category:} & Performance / Memory \\
  \end{tabular}

  \textbf{Description:}
  \texttt{createDeviceLocalBufferExclusive()} uses \texttt{vkAllocateMemory}
  directly instead of going through VMA. This creates a dedicated
  \texttt{VkDeviceMemory} allocation per buffer, causing:

  \begin{itemize}[noitemsep]
    \item Higher memory overhead (minimum allocation granularity).
    \item Fragmentation in device memory.
    \item No suballocation — each buffer gets its own 64+~MB block on some hardware.
  \end{itemize}

  \textbf{Fix:}
  Replace with \texttt{vmaCreateBuffer()} using
  \texttt{VMA\_ALLOCATION\_CREATE\_STRATEGY\_MIN\_MEMORY\_BIT} for large
  allocations, or migrate to VMA suballocation in an existing pool.

\item[\textbf{[P3]}] \textbf{runSingleTimeCommands — Per-Call Alloc Overhead}
  \hfill \severitybox{mediumcolor}{MEDIUM}

  \begin{tabular}{ll}
    \textbf{File:} & \texttt{vulkan/VulkanApp.cpp} \\
    \textbf{Line:} & 1427--1462 \\
    \textbf{Category:} & Performance / CPU \\
  \end{tabular}

  \textbf{Description:}
  Both \texttt{runSingleTimeCommands} and \texttt{runSingleTimeCommandsAsync}
  allocate and free command buffers on every invocation. While
  \texttt{VK\_COMMAND\_POOL\_CREATE\_TRANSIENT\_BIT} mitigates the driver
  overhead, this pattern is called for every staging upload and can cause
  measurable CPU overhead during level loading.

  \textbf{Fix:}
  Use the same ring-buffer pattern already employed for async work:
  pre-allocate 2--4 command buffers per pool and cycle through them.
  Add a lightweight fence wait per slot.

\item[\textbf{[P4]}] \textbf{VegetationRenderer — Map/Unmap Per Operation}
  \hfill \severitybox{mediumcolor}{MEDIUM}

  \begin{tabular}{ll}
    \textbf{File:} & \texttt{vulkan/renderer/VegetationRenderer.cpp} \\
    \textbf{Line:} & 1657--1661 \\
    \textbf{Category:} & Performance / CPU \\
  \end{tabular}

  \textbf{Description:}
  For each chunk generation, the indirect buffer is mapped, written, and
  unmapped. The buffer is host-visible + host-coherent, so the mapping can
  be persistent. Map/Unmap calls are not free — they may involve kernel
  transitions or TLB flushes.

  \textbf{Fix:}
  Map the buffer once at allocation time and keep the pointer. Unmap only
  at buffer destruction. This is safe for host-coherent memory.

\item[\textbf{[S3]}] \textbf{IndirectRenderer::readVisibleCount — Full GPU Drain}
  \hfill \severitybox{mediumcolor}{MEDIUM}

  \begin{tabular}{ll}
    \textbf{File:} & \texttt{vulkan/renderer/IndirectRenderer.cpp} \\
    \textbf{Line:} & 1308 \\
    \textbf{Category:} & Synchronization / Stall \\
  \end{tabular}

  \textbf{Description:}
  \texttt{readVisibleCount()} calls \texttt{app->deviceWaitIdle()} to read
  the visible count back from the GPU. This is a debug/stats function, but
  it introduces a full GPU drain that stalls all queues. On a 16~ms frame,
  this alone can add 10+~ms of latency.

  \textbf{Fix:}
  Use a dedicated \texttt{VkFence} per query. Submit with the fence,
  wait on that fence (with a short timeout), then read back. This avoids
  stalling unrelated work.

\item[\textbf{[C1]}] \textbf{Legacy Render Pass Escape Hatch Present}
  \hfill \severitybox{mediumcolor}{MEDIUM}

  \begin{tabular}{ll}
    \textbf{File:} & \texttt{vulkan/VulkanApp.cpp:4213} \\
    & \texttt{utils/RendererUtils.hpp:99} \\
    \textbf{Category:} & Compliance / Regression Risk \\
  \end{tabular}

  \textbf{Description:}
  Both \texttt{createGraphicsPipeline} and \texttt{buildFullscreenPipeline}
  accept an optional \texttt{VkRenderPass legacyRenderPass} parameter.
  When non-null, they fall back to traditional render pass pipelines,
  bypassing dynamic rendering entirely. This creates a regression vector:
  if any future caller passes a non-null handle, dynamic rendering is
  silently disabled for that pipeline.

  \textbf{Fix:}
  Remove the \texttt{legacyRenderPass} parameter, or mark it
  \texttt{= delete} and refactor any internal use. If backward
  compatibility is needed, add a compile-time flag.

\item[\textbf{[S4]}] \textbf{Swapchain Barrier Uses Conflicting Stage/Access Mask}
  \hfill \severitybox{mediumcolor}{MEDIUM}

  \begin{tabular}{ll}
    \textbf{File:} & \texttt{vulkan/VulkanApp.cpp} \\
    \textbf{Line:} & 4869 \\
    \textbf{Category:} & Synchronization / Correctness \\
  \end{tabular}

  \textbf{Description:}
  The swapchain image transition to \texttt{COLOR\_ATTACHMENT\_OPTIMAL} uses
  \texttt{srcStageMask = COLOR\_ATTACHMENT\_OUTPUT\_BIT} combined with
  \texttt{srcAccessMask = 0}. A non-zero stage with a zero access mask
  is contradictory: the stage says ``make color attachment output visible''
  but the access mask says ``no accesses to make visible.'' While the
  contents are undefined (oldLayout is UNDEFINED), the stage mask should
  be \texttt{TOP\_OF\_PIPE\_BIT} to properly indicate ``nothing to wait on.''

  \textbf{Fix:}
  Change \texttt{srcStageMask} to \texttt{VK\_PIPELINE\_STAGE\_2\_TOP\_OF\_PIPE\_BIT}.

\end{enumerate}

\subsection{Low Severity}

\begin{enumerate}[leftmargin=*]

\item[\textbf{[P5]}] \textbf{Per-Frame Descriptor Writes in SceneRenderer}
  \hfill \severitybox{lowcolor}{LOW}

  \begin{tabular}{ll}
    \textbf{File:} & \texttt{vulkan/renderer/SceneRenderer.cpp} \\
    \textbf{Line:} & 764--773 \\
    \textbf{Category:} & Performance / CPU \\
  \end{tabular}

  \textbf{Description:}
  Descriptor sets for uniform buffers and image samplers are updated every
  frame via \texttt{vkUpdateDescriptorSets}. For resources that change
  infrequently (static textures, constant buffers), this generates
  unnecessary CPU work and driver validation overhead.

  \textbf{Fix:}
  Separate descriptor sets into ``static'' (updated once at init) and
  ``dynamic'' (updated per frame). Use push descriptors or dynamic
  uniform buffers for per-frame data.

\item[\textbf{[B6]}] \textbf{Resource Manager Tracks VkRenderPass Unnecessarily}
  \hfill \severitybox{lowcolor}{LOW}

  \begin{tabular}{ll}
    \textbf{File:} & \texttt{vulkan/VulkanResourceManager.cpp} \\
    \textbf{Line:} & 176, 275, 293, 488 \\
    \textbf{Category:} & Bug / Dead Code \\
  \end{tabular}

  \textbf{Description:}
  \texttt{VulkanResourceManager} tracks \texttt{VkRenderPass} handles in
  its resource tables, but no render passes are ever created (the engine
  uses dynamic rendering exclusively). The tracking code is dead and
  wastes a small amount of memory.

  \textbf{Fix:}
  Remove \texttt{VkRenderPass} tracking from \texttt{VulkanResourceManager}.
  If render passes are needed later, they can be re-added.

\item[\textbf{[C2]}] \textbf{Binary Semaphores Used Alongside Timeline}
  \hfill \severitybox{lowcolor}{LOW}

  \begin{tabular}{ll}
    \textbf{File:} & \texttt{vulkan/VulkanApp.cpp} \\
    \textbf{Line:} & 4799 \\
    \textbf{Category:} & Compliance \\
  \end{tabular}

  \textbf{Description:}
  The present submit signals both a binary semaphore
  (\texttt{renderFinishedSemaphores}) and the frame timeline semaphore.
  While required for WSI compatibility (binary semaphores are needed by
  \texttt{vkQueuePresentKHR}), this creates a dual-signal path that
  complicates synchronization reasoning.

  \textbf{Fix:}
  Acceptable as-is (WSI requires binary semaphores). Document the dual
  signal clearly with a comment at the submission site.

\item[\textbf{[C3]}] \textbf{SPIR-V Targets vulkan1.1 Instead of 1.6}
  \hfill \severitybox{lowcolor}{LOW}

  \begin{tabular}{ll}
    \textbf{File:} & \texttt{Makefile} \\
    \textbf{Line:} & shader compilation rules \\
    \textbf{Category:} & Compliance / Best Practice \\
  \end{tabular}

  \textbf{Description:}
  Shaders are compiled targeting \texttt{vulkan1.1} (SPIR-V 1.3).
  The engine uses Vulkan 1.3/1.4 features. Targeting a higher SPIR-V
  version (1.6 for Vulkan 1.3, or at least 1.5) enables advanced
  features like 64-bit integers, 8/16-bit types, and scalar block layout
  in shaders.

  \textbf{Fix:}
  Change \texttt{--target-env vulkan1.1} to \texttt{vulkan1.3} (or
  \texttt{vulkan1.4} when supported) in the \texttt{glslc} flags.

\end{enumerate}

\newpage

\section{Comparison with Previous Report}

The previous report (perf\_report\_05) flagged 18 issues, of which 4 were resolved,
3 were mitigated, and 11 remain open. The current report adds 8 newly discovered
issues and reaffirms 15 previously identified but unresolved items.

\begin{tabular}{lccc}
\toprule
\textbf{Category} & \textbf{Previous} & \textbf{New This Report} & \textbf{Total Open} \\
\midrule
Bugs & 3 & 4 & 7 \\
Performance & 5 & 2 & 7 \\
Compliance & 2 & 3 & 5 \\
Synchronization & 2 & 5 & 7 \\
\midrule
\textbf{Total} & \textbf{12} & \textbf{14} & \textbf{26} \\
\bottomrule
\end{tabular}

\newpage

\section{Overall Classification}

The engine demonstrates a solid architectural foundation with correct adoption of
modern Vulkan idioms: Synchronization2, dynamic rendering, timeline semaphores,
and VMA-based memory management are all used pervasively and correctly.

However, the following issues prevent the engine from achieving production-grade
stability and performance:

\begin{enumerate}
  \item \textbf{The async command buffer ring (B1) is a latent memory corruption bug}
        that will manifest under sustained heavy load. This is the single highest-priority item.
  \item \textbf{Full pipeline flushes (S2) and GPU-wide drains (S3) nullify the benefits}
        of explicit synchronization and directly cap frame rate.
  \item \textbf{Legacy allocation patterns} (\texttt{vkAllocateMemory}, heap-allocated
        descriptor structs, per-call map/unmap) add unnecessary CPU overhead.
  \item \textbf{Exception safety gaps} (B2, B3, B5) cause resource leaks on error paths,
        making debugging device-lost errors significantly harder.
\end{enumerate}

\subsection*{Overall Severity}

\medskip
\noindent\severitybox{highcolor}{MODERATE}

\begin{quote}
  The engine is functional for development and prototyping but contains
  \textbf{2 critical, 5 high, and 8 medium severity issues} that should
  be addressed before any production or public release. The core architecture
  (dynamic rendering, Synchronization2, timeline semaphores, VMA) is sound
  and follows modern best practices.
\end{quote}

\subsection*{Remediation Priority}

\begin{enumerate}[noitemsep]
  \item[\textbf{P0}] \textbf{Immediate:} B1 (async ring wrap) — corruption risk.
  \item[\textbf{P1}] \textbf{This week:} S2 (ALL\_COMMANDS flush), S3 (deviceWaitIdle), S1 (shadow barrier).
  \item[\textbf{P2}] \textbf{This sprint:} B2/B3/B5 (exception safety), P1 (staging ring double-free), P2\\(direct vkAllocateMemory).
  \item[\textbf{P3}] \textbf{Backlog:} P3 (runSingleTimeCommands), P4 (map/unmap), C1 (legacy escape hatch), C3 (SPIR-V target).
\end{enumerate}

\newpage

\section{Fix Instructions Summary}

\begin{enumerate}[noitemsep]
  \item B1 — Add fence-per-slot to async ring; wait before reuse.
  \item B2 — Call cleanup() before throwing in device-lost handler.
  \item B3 — Use RAII wrappers for buffer allocations in async upload.
  \item B4 — Set acquire timeout to 1~s; handle VK\_TIMEOUT gracefully.
  \item B5 — Replace new/delete with stack allocation for descriptor info structs.
  \item B6 — Remove VkRenderPass tracking from VulkanResourceManager.
  \item P1 — Add double-release guard to StagingRingBuffer::release().
  \item P2 — Migrate createDeviceLocalBufferExclusive to VMA.
  \item P3 — Pre-allocate command buffer ring for single-time commands.
  \item P4 — Persistent map host-coherent buffers in VegetationRenderer.
  \item P5 — Split descriptor updates into static/dynamic sets.
  \item C1 — Remove legacyRenderPass parameter from pipeline creation.
  \item C2 — Add explanatory comment at the dual-signal submit site.
  \item C3 — Bump SPIR-V target to vulkan1.3+ in glslc flags.
  \item S1 — Fix shadow cascade barrier to wait on LATE\_FRAGMENT\_TESTS.
  \item S2 — Replace ALL\_COMMANDS\_BIT with targeted stage masks.
  \item S3 — Use per-query fence in readVisibleCount instead of deviceWaitIdle.
  \item S4 — Use TOP\_OF\_PIPE\_BIT for swapchain acquire barrier.
\end{enumerate}

\end{document}
